% !TeX program = xelatex

\documentclass{article}
\usepackage{fontspec}
%% https://www.ctan.org/pkg/polyglossia
\usepackage{polyglossia}
\setdefaultlanguage{chinese}
\setotherlanguage{english}
% Package noto, notoCJKsc
\setmainfont{Sarasa Mono SC}[Script=CJK]
\newfontfamily\chinesefont{Sarasa Mono SC}[Script=CJK]
% Package arphic-ttf
%\setmainfont{bkai00mp.ttf}[Script=CJK]% ZenKai-Medium
%\setmainfont{bsmi00lp.ttf}[Script=CJK]% ShanHeiSun-Light
%\setmainfont{gbsn00lp.ttf}[Script=CJK]% BousungEG-Light-GB
%\setmainfont{gkai00mp.ttf}[Script=CJK]% GBZenKai-Medium
\begin{document}
\parindent0em

\begin{center}
  \abstractname
\end{center}
\begin{english}
  This package provides an alternative to Babel for users of XeLaTeX
  and LuaLaTeX. This version includes support for over 80 different
  languages, some of which in different regional or national
  varieties, or using a different writing system.
  \footnote{%
  This is a footnote.}

  This is today: \today
\end{english}

\section{(简体中文, Simplified Chinese)}

该软件包为使用XeLaTeX和LuaLaTeX的用户提供了Babel的替代方案。此版本支持超过80种不同的语言，其中一些具有不同的区域或国家变体，或使用不同的书写系统。

\today

localnumeral: \localnumeral{1863}, chinesenumeral: \chinesenumeral{1863}

\renewfontfamily\chinesefont{Sarasa Mono SC}[Script=CJK]
\begin{chinese}[variant=traditional,numerals=chinese]

  \section{(繁體中文, Traditional Chinese)}

  該軟體包為使用XeLaTeX和LuaLaTeX的用戶提供了Babel的替代方案。此版本支持超過80種不同的語言，其中一些具有不同的區域或國家變體，或使用不同的書寫系統。

  \today

  localnumeral: \localnumeral{1863}, chinesenumeral: \chinesenumeral{1863}

\end{chinese}

\section{numeral}
\begin{english}

  numerals=arabic: \textchinese[numerals=arabic]{\localnumeral{1863},
  \chinesenumeral{1863}}

  numerals=chinese:
  \textchinese[numerals=chinese]{\localnumeral{1863}, \chinesenumeral{1863}}

\end{english}

\end{document}
